\paragraph{Block periods of false lifts}\label{par:false-lift-periods}
Let \(W,V\) be the two spaces of the descent pair, of ranks \(r_d,n\).
Take \(Y\) to be the conjugate of a matrix coefficient of the form being
tested. For the highest weight coefficient, the descent with the harmonic
polynomial and the finite block datum vanishes. With the finite datum fixed,
replacing the archimedean polynomial by any Fock polynomial gives the
convolution of this zero function with a test function on the real points
of \(U(W)\). These descents therefore vanish as well. A change of the
character multiplies the descent by a factor on the determinant of its
argument, so the conclusion holds with the normalization used here.
\leanlinks{\leansrc{New/OneLift/FalseLifts/PeriodsOfFalseLifts}{From the harmonic datum to arbitrary polynomials}}

At the places \(w\) of the level, decompose a finite test function into
indicators of sufficiently small residue classes. Write the center as
\(z_w=\binom{z'_w}{z''_w}\). Suppose that all selected row blocks \(z'_w\)
are invertible and \(z''_w(z'_w)^{-1}\) is integral, with zero reduction at
the chosen places. Right translation by \((z'_w)^{-1}\) puts the center
where the block datum has value \(1\), and that datum is constant on a
sufficiently small class. The character-weighted average of the indicator
over level translations is a scalar multiple of the translated block
datum. The level group acts transitively on its support, and the
transformation rule determines values from one value. Pairing with the
descent is invariant under the averaging, so this pairing is zero.

If the condition fails at a place, use the kernel of the selected row block,
or clear the denominator of a row of \(z''_w(z'_w)^{-1}\) to obtain a
primitive row. This gives an integral row \(l\) with \(lz_w=0\) whose
reduction \(\lambda\) on the selected rows is nonzero. At a twist place
the remaining entries of \(l\) can be taken to have zero reduction. Lift
each invertible \(s=1+v\lambda\) over the residue field to an integral
matrix \(1+\widetilde v l\). It fixes \(z_w\) and preserves the residue
class. On the form it acts as \(\chi_w(\det s)^{-1}\bar s\), where
\(\bar s\) is the block action of \(s\) at a chosen place and \(1\) at a
twist place. Adding the descents after these translations gives the
original descent value times the order of the group. At chosen places
the sum is zero by the hyperplane condition. At twist places the
determinant \(1+\lambda v\) runs with constant multiplicity over the
units, so the sum is zero by \(\chi_w\ne1\). This proves vanishing of the
descents to \(W\) for all data that are vacuum off the level.
\leanlinks{\leansrc{New/OneLift/FalseLifts/BottomOfTower}{Integrals on small classes and hyperplane sums}}

Add the hyperbolic planes one at a time. After adding \(i\ge1\) planes,
put \(m'=r_d+2i\), \(k'=n-m'\), and write \(F\) for the descent.
In the mixed model the zero term along the last isotropic line integrates
to zero by the induction hypothesis for \(i-1\) planes. The remaining
kernel is uniformly bounded. Indeed, the definite target has height
\(1\). If the source height \(H\) exceeds \(2\), rational translation
makes the isotropic line have length at most \(2/H\). Subtracting the zero
term gives a functional that kills its translation average. Decay of
arbitrary order under the resulting short translations absorbs the moderate
growth \(O(H^A)\). For \(H\le2\), moderate growth suffices. Integrating
against the form on the compact target quotient makes \(F\) bounded.

Let \(q\) be the residue cardinality at the split place off the level and
\(P_{d,w}(T)\) the rank-\(d\) Hecke polynomial. The original form is killed
by \(P_{n,w}(X)\) with \(|X|=q^{-(n-1-r_d/2)}\).
Transferring the local Hecke action on the vacuum makes \(F\) spherical
and gives
\[
 \prod_{j=0}^{k'-1}(1-q^jX_0)\,P_{m',w}(X')F=0,\qquad
 X_0=\chi(\mathrm{Frob}_{\bar w})q^{m'/2}X,\qquad
 X'=\frac{q^{k'}X_0}{q^{n/2}} .
\]
Here \(\chi\) is the base change of the descent character.
The inequalities \(i\ge1\) and \(2r_d+2i+3<n\) imply
\(q^{k'+m'/2}|X|\le1\) and \(q^{m'+1+k'/2}|X|\le1\).
The first makes the product nonzero and the second gives
\(q^{m'+1}|X'|\le1\). In the expansion of \(P_{m',w}\) over subspaces of
the residue vector space, the zero subspace gives the identity and the
Gaussian binomial formula gives the sum of the absolute coefficients
\[
 \prod_{j=0}^{m'-1}(1+q^j|X'|)
 \le\exp\!\left(\frac{1}{q(q-1)}\right)<2 .
\]
Taking the supremum after isolating the identity term bounds
\(\|F\|_\infty\) by a multiple strictly less than \(1\) of itself.
Thus \(F=0\). Translates of the highest vector span the coefficient space,
so the result holds for every matrix coefficient.
\leanlinks{\leansrc{New/OneLift/FalseLifts/HyperbolicPlaneStep}{Subtracting the zero term and using a small Hecke root}}

Write \(V=V_1\oplus V_2\) with \(\dim V_1=N_d\) and \(2r_d+2N_d+3<n\).
Complete the kernel of \(W\) on \(V_2\) occurring in the period by a kernel
on \(V_1\), with inverse block characters. The \(V_1\) datum can be chosen
so that its kernel is nonzero at the fixed arguments. Multiply \(Y\) by
the resulting kernel on \(U(V)\) to get \(Y_2\). Multiplication of kernels
in the source makes the descent of \(Y_2\) to the hyperbolic space a
descent of \(Y\) to \(W\oplus\mathbb H^{N_d}\), which vanishes.
Finite right translations supported at the level preserve the condition
on the data.

Let \(I(g)\) be the translated block period of \(Y_2\) over \(U(V_2)\),
and \(x(g)\) its corresponding frame of \(V_1\). The Fourier coefficient
in the mixed model along the isotropic frame, for a Heisenberg translate
of the vacuum, introduces the weight
\(\psi_u(x(g))\,|e(\alpha_0x(g))|^2\).
Here \(\psi_u\) is the additive character pairing the frame entries
with \(u\), omitting the components held fixed by the vacuum condition.
Integrating first over \(U(V_2)\) gives zero for the integral of this weight
times \(I(g)\), for every \(u\). On the vacuum support, changes in the
omitted components can be absorbed by the integral unitary level away
from the selected places. The conjugate of the period thus descends
continuously through the values of the additive characters. Applying
Stone--Weierstrass to the algebra generated by their products and
conjugates gives
\[
 \int |e(\alpha_0x(g))|^2\,|I(g)|^2\,\mathrm dg=0 .
\]
The integrand is continuous and nonnegative. The vacuum weight is positive
at the frame giving the original period, so that period is zero. Divide
by the nonzero first-block factor to obtain vanishing of the desired
\(V_2\) period. On level classes the integral is a sum weighted by
stabilizer orders.
\leanlinks{\leansrc{Lib/Weil/Periods/DescentVanishing}{Recovering the period from Fourier coefficients of the frame}
  \leanlinksep \leansrc{New/OneLift/FalseLifts/PeriodsOfFalseLifts}{Vanishing of the class sum on the block}}
